\cotefile{156-5}
\scanpage{42}\mspage{41}
c) L'\unsure{ens|ens.} des fermés pour ${\unsure{\mathcal{T}|\tau}}'$ est en corr. \unsure{biuniv.} avec l'\unsure{ens|ens.} des couples
\[
\begin{aligned}
(B,A_{\unsure{u|n|U}}) \qquad & B\unsure{\subset|\in}T\ \text{fermé pour }\mathcal{T}'_0\\
& A_u\ \text{fermé de }\unsure{U|u|n}=S\setminus T\ \text{pour }\mathcal{T}\mid(S\setminus T)
\end{aligned}
\]
\[
\text{tels que}\qquad \overline{A_u}\cap T\subset B \qquad
\text{(adh. prise pour }\mathcal{T}\text{)}\ \deleted{\unsure{\mathcal{T})}}.
\]

d) Si $\mathcal{T}'_0$ et $\mathcal{T}$ \unsure{ultra topologiques|ultratopologiques}, $\mathcal{T}'$ aussi\unsure{.} La relation de préordre associée \deleted{\unsure{se définit}} \inserted{se décrit} ainsi. Pour $x,y\in S$, on a $x\underset{\mathcal{T}'}{\leq}y$ sss
\[
\left\{\begin{aligned}
&\text{ou bien}\quad x,y\in T,\qquad x\underset{\mathcal{T}'}{\leq}y\\
&\left.\begin{aligned}
&\text{ou bien}\quad x,y\in U=S\setminus T,\qquad x\underset{\mathcal{T}}{\leq}y\\
&\text{ou bien}\quad x\in T,\ y\in U,\qquad x\underset{\mathcal{T}}{\leq}y
\end{aligned}\right]
\end{aligned}\right.
\]
\margin{en d'autres termes, si $y\in U$, on a $x\underset{\mathcal{T}'}{\leq}y$ sss \unsure{a} $x\underset{\mathcal{T}}{\leq}y$.}

e) \inserted{Sous les conditions de d)} Considérons les ens ordonnés $I,I'$ associés à $\mathcal{T},\mathcal{T}'$, et la partie \inserted{fermée} $J\subset I$ correspondant à la partie fermée $T$ de $(S,\mathcal{T})$, enfin l'application $\varphi:I'\unsure{\mathbin{\text{---}}|\longrightarrow}I$ canonique.
\diagram{\begin{tikzcd}[ampersand replacement=\&,column sep=small,row sep=large]
\varphi^{-1}(J)={} \& J' \arrow[d] \& \subset \& I' \arrow[d,"\varphi"] \& \overline{\mathcal{T}'} \arrow[d] \\
{} \& J \& \subset \& I \& \overline{\mathcal{T}}
\end{tikzcd}}

\deleted{Alors}, \deleted{\unsure{la|le}} Soit $\overline{\mathcal{T}'}$ \deleted{est} la \unsure{top|top.} de $\mathcal{T}'$ déduite de $\mathcal{T}'$ par passage au quotient \unsure{(ou encore,} associée à la relation \uline{d'ordre} sur $I'$ associé au préordre), $\widetilde{\mathcal{T}}$ la topologie sur $I'$ image inverse de $\overline{\mathcal{T}}$. Alors $\overline{\mathcal{T}'}$ est la moins fine des top. sur $I'$ qui \inserted{\unsure{sont plus fines que} $\widetilde{\mathcal{T}}$ et} induisent sur $J'$ la $\unsure{=}$ \deleted{\unsure{top}} topologie $\overline{\mathcal{T}}\mid J'$. Elle est
